% ============================================================
% Chapter 6: Conclusion and Future Work
% ============================================================
\chapter{Conclusion and Future Work}
\label{ch:conclusion}

This chapter summarizes the key findings of the thesis (\Cref{sec:summary}), discusses limitations (\Cref{sec:limitations}), and outlines directions for future research (\Cref{sec:future_work}).

% ============================================================
\section{Summary of Contributions and Findings}
\label{sec:summary}

This thesis presented SustainRL-Bench, a comprehensive benchmarking framework for evaluating reinforcement learning algorithms across three sustainable energy environments---electric vehicle charging coordination, building HVAC control, and cogeneration plant dispatch---along three experimental axes: perturbation robustness, safety constraints, and multi-agent coordination. We summarize the principal findings below, organized by research question.

\subsection{Perturbation Robustness (RQ1, RQ2, RQ5)}

The systematic evaluation of PPO, SAC, and TD3 under state, action, and dynamics perturbation across ten or more perturbation levels per environment yielded three key insights:

\begin{enumerate}
    \item \textbf{Environment structure dominates perturbation sensitivity.} EV Charging policies are exceptionally robust to all perturbation types ($<1\%$ degradation at moderate perturbation levels), owing to the reward's dependence on aggregate daily charging patterns rather than instantaneous observations. In contrast, Building policies degrade rapidly ($>100\%$ degradation at $\sigma \geq 0.15$) because the RC thermal model amplifies per-timestep errors through coupled zone dynamics. Cogeneration exhibits algorithm-dependent sensitivity, with PPO virtually immune and SAC degrading catastrophically under state perturbation.

    \item \textbf{State perturbation is far more damaging than action perturbation.} Across all three environments and all algorithms, state perturbation causes significantly greater performance degradation than action perturbation at equivalent scales. This finding has a clear practical implication: for real-world deployment of RL-based energy controllers, investment in sensor accuracy yields greater returns than investment in actuator precision.

    \item \textbf{PPO is the most perturbation-robust algorithm.} PPO consistently outperforms SAC and TD3 in perturbation robustness, with the most dramatic advantage in Cogeneration (PPO: $<0.3\%$ degradation; SAC: $>2000\%$). PPO's on-policy updates and clipped surrogate objective produce naturally conservative policies that tolerate input perturbations, whereas SAC's entropy-maximizing exploration can amplify perturbation through the policy output.
\end{enumerate}

\subsection{Safe Reinforcement Learning (RQ3)}

The formulation of each environment as a Constrained MDP with environment-specific cost signals revealed both successes and fundamental limitations:

\begin{enumerate}
    \item \textbf{Reward--cost orthogonality is necessary for effective CMDP constraints.} When the cost signal measures an objective already well-optimized by the reward (as in Building, where the reward's comfort term with $\beta = 0.5$ already minimizes temperature error), all cost limits produce identical behavior. In contrast, when the cost is genuinely orthogonal to the reward (Cogeneration: turbine ramping stress; EV Charging: excess charge violations), CMDP algorithms achieve precise constraint tracking with clear monotonic differentiation across cost limits.

    \item \textbf{PPOLag provides the most reliable constraint satisfaction.} Across all environments, PPOLag achieves the best cost limit differentiation and precise constraint tracking. Its Lagrangian relaxation degrades gracefully under inactive constraints ($\lambda \to 0$ reduces to standard PPO), unlike trust-region methods (CPO, OnCRPO) that can become numerically unstable when constraints are nearly satisfied.

    \item \textbf{Off-policy CMDP methods fail.} SACLag consistently failed to satisfy constraints across all three environments due to fundamental incompatibility between replay buffer staleness and Lagrange multiplier updates. This negative finding demonstrates that on-policy methods are currently necessary for reliable CMDP-based safe RL in these domains.
\end{enumerate}

\subsection{Multi-Agent Reinforcement Learning (RQ4)}

We designed and implemented multi-agent versions of all three environments using the PettingZoo parallel interface with physically motivated agent decompositions: per-station agents in EV Charging (54 agents), per-zone agents in Building HVAC, and per-turbine agents in Cogeneration (4 agents). Training with Ray RLlib using PPO, SAC, APPO, and IMPALA is ongoing; preliminary results and the complete single-agent versus multi-agent comparison will be reported in the final version of this thesis.

\subsection{Cross-Environment Insights (RQ5)}

The most consistent cross-cutting finding is that \emph{environment structure matters more than algorithm choice}: the physical characteristics of each energy system---reward aggregation timescale, dynamics coupling, and reward--cost alignment---determine the effectiveness of perturbation robustness and safety interventions more than the choice of RL algorithm or hyperparameters. This suggests that practitioners should invest primarily in understanding their specific energy system's characteristics before selecting algorithms or robustness strategies.

% ============================================================
\section{Limitations}
\label{sec:limitations}

We acknowledge the following limitations of the present work:

\begin{enumerate}
    \item \textbf{Single-seed training for perturbation experiments.} While evaluation uses 100 episodes per configuration with bootstrap confidence intervals, each perturbation-level configuration was trained with a single seed due to computational constraints. Future work should extend to multiple training seeds to capture inter-seed variability in the learned policies themselves.

    \item \textbf{Simulation-only evaluation.} All experiments are conducted in simulation environments. While SustainGym environments are grounded in real-world data and physics-based models, the sim-to-real transfer gap remains unquantified. Real-world deployment would introduce additional challenges such as non-stationary dynamics, hardware failures, and communication delays.

    \item \textbf{Building Safe RL negative result.} The Building environment's reward--cost alignment prevents meaningful CMDP constraint differentiation with the tested cost functions. While we report this as a principled negative finding (illustrating the orthogonality requirement), it means that one of the three environments does not benefit from the Safe RL axis. Alternative cost formulations---such as peak power consumption or zone-specific comfort constraints---may yield better results.

    \item \textbf{Limited MARL algorithm coverage.} The MARL experiments use independent learning with parameter sharing rather than more sophisticated CTDE algorithms such as QMIX or MADDPG. While independent learning provides a strong baseline, the comparison does not capture the full potential of cooperative MARL methods.

    \item \textbf{Fixed environment configurations.} Each environment uses a single configuration (e.g., OfficeSmall/Hot\_Dry for Building, Caltech site for EV Charging). Generalization across building types, weather profiles, or charging network topologies is not evaluated.

    \item \textbf{No combined perturbation evaluation.} Perturbation channels are evaluated independently. Real-world systems experience simultaneous state, action, and dynamics perturbation, and the interaction effects between channels may be non-additive.
\end{enumerate}

% ============================================================
\section{Future Work}
\label{sec:future_work}

Several promising directions emerge from the findings and limitations of this thesis:

\subsection{Perturbation-Aware Training and Curriculum Learning}

Our results show that policies trained under clean conditions can degrade severely under deployment perturbation (particularly in Building). A natural extension is \emph{curriculum-based perturbation training}, where the perturbation level is gradually increased during training to produce policies that are robust without sacrificing clean-condition performance. Domain randomization across perturbation levels and adversarial perturbation training~\citep{pinto2017rarl} are complementary approaches worth investigating in the energy domain.

\subsection{Cross-Environment Transfer Learning}

The three SustainGym environments share structural similarities (episodic 24-hour horizons, multi-component control, energy-comfort/cost tradeoffs) that may enable transfer learning. Pre-training on one environment and fine-tuning on another could reduce sample complexity, particularly for the data-hungry Building and Cogeneration environments. The shared perturbation injection framework provides a natural basis for evaluating whether perturbation robustness transfers across environments.

\subsection{Improved Safe RL Formulations}

The Building Safe RL negative result motivates exploration of alternative cost formulations that are genuinely orthogonal to the reward. Promising candidates include:
\begin{itemize}
    \item \textbf{Peak power constraints}: limiting the maximum instantaneous HVAC power draw, which is orthogonal to the time-averaged energy penalty in the reward.
    \item \textbf{Zone-specific comfort bounds}: enforcing hard temperature limits on individual zones rather than penalizing average deviation.
    \item \textbf{Multi-constraint CMDPs}: simultaneously constraining multiple cost signals (e.g., both energy and comfort in Building, both ramping and emissions in Cogeneration).
\end{itemize}

Additionally, the SACLag failure motivates research into off-policy CMDP methods that address replay buffer staleness, such as importance-weighted Lagrangian updates or cost-conditioned replay buffers.

\subsection{Robustness of Constrained RL Under Perturbation}

A critical open question is whether CMDP-based safety guarantees remain valid under the realistic perturbation conditions evaluated in \Cref{ch:robustness}. This thesis evaluates perturbation robustness and safe RL as independent axes; however, real-world energy systems will experience both simultaneously---a building HVAC controller must satisfy comfort constraints \emph{while} operating with noisy temperature sensors, and a cogeneration dispatcher must limit turbine ramping stress \emph{despite} uncertain ambient measurements.

Recent work by \citet{gu2025robustgymnasium} demonstrates that this concern is well-founded. Their Robust-Gymnasium benchmark evaluates safe RL algorithms (PCRPO and CRPO) on SafetyWalker2d under action perturbation and cost signal perturbation, finding that CRPO's constraint satisfaction degrades rapidly under both perturbation types, while PCRPO exhibits greater resilience---and in some cases, perturbation during training actually \emph{improves} deployment performance through a regularization effect. However, their evaluation is limited to two algorithms, a single locomotion environment, and a single cost limit, leaving the interaction between perturbation robustness and constraint satisfaction largely unexplored in applied domains.

Our SustainRL-Bench infrastructure is uniquely positioned to investigate this interaction comprehensively. The three independent perturbation channels (state, action, and dynamics) and four CMDP algorithms (PPOLag, CPO, OnCRPO, FOCOPS) are already implemented and validated independently; combining them requires no additional environment engineering. Several specific research questions arise:

\begin{itemize}
    \item \textbf{Constraint violation under perturbation}: Do policies trained with strict cost limits (e.g., $c_{\text{limit}} = 1$ for EV Charging) maintain constraint satisfaction when deployed under state perturbation? The Lagrange multiplier in PPOLag is calibrated to the training distribution; distribution shift from perturbation may cause the learned $\lambda$ to under- or over-correct.

    \item \textbf{Perturbation-aware CMDP training}: Training CMDP algorithms under perturbation simultaneously could yield policies that satisfy safety constraints robustly. This would combine the curriculum-based perturbation training discussed in \Cref{sec:future_work} with constrained optimization, effectively solving a \emph{robust} CMDP.

    \item \textbf{Environment-dependent interaction effects}: Our robustness results show that EV Charging is naturally perturbation-tolerant while Building is highly sensitive. This asymmetry likely extends to the safe RL setting: Building's tight thermal coupling may cause perturbation to amplify cost violations, while EV Charging's reward aggregation may buffer constraint satisfaction against moderate perturbation.

    \item \textbf{Cost signal corruption}: Following \citet{gu2025robustgymnasium}'s cost signal attack, perturbation applied directly to the observed cost (rather than to states or actions) tests a fundamentally different failure mode---the agent receives incorrect safety feedback, potentially learning to violate constraints while believing they are satisfied.
\end{itemize}

This direction represents a natural intersection of the perturbation robustness and safe RL axes developed in this thesis, and would provide actionable guidance for deploying constrained RL controllers in real-world energy systems where both safety and robustness are non-negotiable requirements.

\subsection{Scalable Multi-Agent Methods}

The EV Charging environment with 54 agents provides a challenging testbed for scalable MARL. Future work should investigate:
\begin{itemize}
    \item \textbf{Communication-constrained MARL}: where agents can only exchange limited information with neighbors, reflecting realistic communication constraints in charging networks.
    \item \textbf{Heterogeneous agents}: different agent types with asymmetric capabilities (e.g., fast-charging vs.\ standard stations).
    \item \textbf{Safe MARL}: combining CMDP constraints with multi-agent coordination, where safety constraints apply to the collective system behavior.
\end{itemize}

\subsection{Real-World Deployment Considerations}

Bridging the sim-to-real gap for energy RL requires addressing several practical challenges:
\begin{itemize}
    \item \textbf{Online adaptation}: policies that can adapt to non-stationary dynamics (seasonal weather changes, evolving occupancy patterns, equipment degradation) without full retraining.
    \item \textbf{Safety during exploration}: deploying RL agents in physical energy systems requires exploration strategies that never violate hard safety constraints, even during the initial learning phase.
    \item \textbf{Interpretability}: energy system operators require understanding of \emph{why} an RL agent makes specific decisions, motivating research into interpretable RL policies for energy applications.
\end{itemize}

\subsection{Extended Benchmark Coverage}

The SustainRL-Bench framework can be extended along several dimensions:
\begin{itemize}
    \item \textbf{Additional environments}: incorporating SustainGym's electricity market and data center cooling environments.
    \item \textbf{Additional algorithms}: evaluating model-based RL methods (e.g., Dreamer, MBPO) that may offer improved sample efficiency and perturbation robustness through learned dynamics models.
    \item \textbf{Combined perturbation evaluation}: systematically evaluating all combinations of state, action, and dynamics perturbation to characterize interaction effects.
    \item \textbf{Longer horizons}: extending the evaluation to multi-day or seasonal horizons to capture long-term dynamics and distribution shift.
\end{itemize}

% ============================================================
\section{Closing Remarks}
\label{sec:closing}

This thesis demonstrates that reinforcement learning holds significant promise for sustainable energy system control, but that principled evaluation across perturbation, safety, and multi-agent dimensions is essential for understanding the practical deployment landscape. The SustainRL-Bench framework, associated environments, and experimental infrastructure provide a foundation for the research community to systematically evaluate and improve RL algorithms for real-world energy applications. We hope that the findings, negative results, and open questions presented in this work will guide future research toward more robust, safe, and scalable RL-based energy controllers.
